\documentclass[10pt,letterpaper]{article}
\usepackage[T1]{fontenc}
\usepackage[utf8]{inputenc}
\usepackage{lmodern}
\usepackage[letterpaper,margin=0.84in,headheight=15pt]{geometry}
\usepackage[table]{xcolor}
\definecolor{DeepNavy}{HTML}{173450}
\definecolor{Teal}{HTML}{087F8C}
\definecolor{Slate}{HTML}{536578}
\definecolor{SoftBlue}{HTML}{F1F6FA}
\definecolor{SoftGray}{HTML}{F6F8FA}
\definecolor{Crimson}{HTML}{B34843}
\usepackage{microtype}
\setlength{\emergencystretch}{1.5em}
\usepackage{booktabs,longtable,tabularx,array}
\usepackage{amssymb}
\usepackage{ragged2e}
\usepackage{enumitem}
\usepackage{fancyhdr,lastpage}
\usepackage{titlesec}
\usepackage{listings}
\usepackage[most]{tcolorbox}
\usepackage{needspace}
\usepackage{hyperref}
\hypersetup{colorlinks=true,linkcolor=DeepNavy,urlcolor=Teal,citecolor=DeepNavy,pdfauthor={Software Engineering Technical Documentation}}
\setlength{\parindent}{0pt}
\setlength{\parskip}{5pt plus 1pt}
\setlist[itemize]{itemsep=2pt,topsep=3pt,leftmargin=1.55em}
\setlist[enumerate]{itemsep=2pt,topsep=3pt,leftmargin=1.7em}
\titleformat{\section}{\Large\bfseries\color{DeepNavy}}{\thesection}{0.75em}{}[\vspace{-3pt}\color{Teal}\titlerule]
\titleformat{\subsection}{\large\bfseries\color{DeepNavy}}{\thesubsection}{0.65em}{}
\titleformat{\subsubsection}{\normalsize\bfseries\color{Teal}}{\thesubsubsection}{0.55em}{}
\pagestyle{fancy}
\fancyhf{}
\fancyhead[L]{\sffamily\footnotesize MISRA C:2023 | Engineering Documentation}
\fancyhead[R]{\sffamily\footnotesize\color{Slate} Independent companion}
\fancyfoot[L]{\sffamily\scriptsize Original explanatory material; consult licensed MISRA standard}
\fancyfoot[R]{\sffamily\scriptsize\thepage\ / \pageref{LastPage}}
\renewcommand{\headrulewidth}{0.3pt}
\newcommand{\ID}[1]{\textsf{\bfseries #1}}
\newcommand{\SourceNote}[1]{\par\smallskip{\footnotesize\color{Slate}\textit{Trace: #1.}}\par}
\newtcolorbox{keybox}[1][]{enhanced,breakable,colback=SoftBlue,colframe=Teal,boxrule=0.5pt,arc=1mm,title={Key point},fonttitle=\bfseries,#1}
\newtcolorbox{riskbox}[1][]{enhanced,breakable,colback=SoftGray,colframe=Crimson,boxrule=0.5pt,arc=1mm,title={Verification boundary},fonttitle=\bfseries,#1}
\lstdefinestyle{engineerc}{language=C,basicstyle=\ttfamily\footnotesize,keywordstyle=\color{DeepNavy}\bfseries,commentstyle=\color{Slate}\itshape,stringstyle=\color{Teal},numbers=left,numberstyle=\tiny\color{Slate},stepnumber=1,numbersep=8pt,breaklines=true,breakatwhitespace=false,frame=single,framesep=5pt,rulecolor=\color{Slate!25},backgroundcolor=\color{SoftGray},tabsize=2,showstringspaces=false,columns=fullflexible,keepspaces=true}
\lstdefinestyle{plaincode}{basicstyle=\ttfamily\footnotesize,numbers=none,breaklines=true,frame=single,rulecolor=\color{Slate!25},backgroundcolor=\color{SoftGray},columns=fullflexible,keepspaces=true}
\newcommand{\TrainingTag}[1]{\textsf{\footnotesize\color{Teal}\bfseries #1}}
\newcommand{\DocNotes}{\begin{keybox}This publication is an independently written educational and operational companion. It is \emph{not} the licensed MISRA guideline text, a MISRA-endorsed tool, or a certification of project compliance. Identifiers and published classifications are used for traceability; check exact requirement wording, amplification, applicability, and exceptions in the licensed standard.\end{keybox}}
\AtBeginDocument{\color{black}}

\title{\vspace{-2.3em}\color{DeepNavy}\Huge\bfseries C Code Pattern Atlas\\[0.45em]\large\normalfont\color{Teal} Problematic and preferred patterns with verification exercises}
\author{Independent Software Engineering Reference Series}
\date{Edition 1.0 \quad | \quad October 2026}
\begin{document}
\maketitle
\vspace{-0.4em}
\begin{center}\color{Slate}\rule{0.83\linewidth}{0.4pt}\end{center}
\DocNotes
\tableofcontents
\newpage

\section{Reading and validating the examples}
Each pair isolates a teaching issue. ``Problematic'' indicates a defect or risk with respect to the cited guideline; it does not assert the snippet violates \emph{only} that guideline. ``Preferred'' illustrates one aspect of a safer pattern but is not a complete MISRA-compliant program. Supporting configuration, declarations, project context, and all other applicable guidelines still require review. Some bad examples intentionally exhibit undefined behavior and must \textbf{never} be executed in production.
\subsection{Unsigned literal intent}
\TrainingTag{Cross-reference: Rule 7.2}
\subsubsection*{Problematic}
\begin{lstlisting}[style=engineerc]
#include <stdint.h>
uint32_t mask = 0x80000000; /* type depends on model */
\end{lstlisting}
\subsubsection*{Preferred pattern}
\begin{lstlisting}[style=engineerc]
#include <stdint.h>
uint32_t mask = UINT32_C(0x80000000);
\end{lstlisting}
\textbf{Engineering reasoning.} Integer literal type depends on suffix and target integer widths. The fixed-width macro makes the intended integer type visible; check that the macro invocation meets its own constraints.

\textbf{Verification exercise.} Compile under the target ABI and inspect the diagnostic mappings for Rule 7.2 and Rule 7.5.

\subsection{Automatic storage initialization}
\TrainingTag{Cross-reference: Rule 9.1}
\subsubsection*{Problematic}
\begin{lstlisting}[style=engineerc]
uint16_t scale(uint16_t input)
{
    uint16_t factor;
    if (input > 10U) { factor = 2U; }
    return (uint16_t)(input * factor); /* may be unset */
}
\end{lstlisting}
\subsubsection*{Preferred pattern}
\begin{lstlisting}[style=engineerc]
uint16_t scale(uint16_t input)
{
    uint16_t factor = 1U;
    if (input > 10U) { factor = 2U; }
    return (uint16_t)(input * factor);
}
\end{lstlisting}
\textbf{Engineering reasoning.} The second version initializes the local on every control-flow path. Product requirements must still bound the multiplication result and justify the narrowing conversion.

\textbf{Verification exercise.} Add zero, 10, 11 and maximum-input tests; statically analyze every path.

\subsection{Narrowing essential types}
\TrainingTag{Cross-reference: Rule 10.3; Rule 10.8}
\subsubsection*{Problematic}
\begin{lstlisting}[style=engineerc]
uint32_t sensor = 70000U;
uint16_t sample = sensor; /* narrowing */
\end{lstlisting}
\subsubsection*{Preferred pattern}
\begin{lstlisting}[style=engineerc]
uint32_t sensor = 70000U;
uint16_t sample = 0U;
if (sensor <= UINT16_MAX)
{
    sample = (uint16_t)sensor;
}
else
{
    /* reject sample and report range error */
}
\end{lstlisting}
\textbf{Engineering reasoning.} An explicit guard plus intentional cast makes the range assumption auditable, but the final policy for rejecting values belongs to the function contract.

\textbf{Verification exercise.} Prove the guard dominates the cast and test both sides of UINT16\_MAX.

\subsection{Pointee qualifiers}
\TrainingTag{Cross-reference: Rule 11.8}
\subsubsection*{Problematic}
\begin{lstlisting}[style=engineerc]
void update(const uint32_t *register_value)
{
    uint32_t *p = (uint32_t *)register_value;
    *p = 3U;
}
\end{lstlisting}
\subsubsection*{Preferred pattern}
\begin{lstlisting}[style=engineerc]
uint32_t read_value(const uint32_t *register_value)
{
    return *register_value; /* read-only contract */
}
\end{lstlisting}
\textbf{Engineering reasoning.} Casting away const erases the caller-visible immutability guarantee. If a device register is writable, describe that capability directly in its interface instead.

\textbf{Verification exercise.} Search for casts that remove const, volatile or atomic qualifiers.

\subsection{Array indexing and pointer bounds}
\TrainingTag{Cross-reference: Rule 18.1; Dir 4.1}
\subsubsection*{Problematic}
\begin{lstlisting}[style=engineerc]
uint8_t read_at(const uint8_t *data, size_t index)
{
    return data[index]; /* caller supplied no length */
}
\end{lstlisting}
\subsubsection*{Preferred pattern}
\begin{lstlisting}[style=engineerc]
bool read_at(const uint8_t *data, size_t n,
             size_t index, uint8_t *out)
{
    bool ok = false;
    if ((data != NULL) && (out != NULL) && (index < n))
    {
        *out = data[index];
        ok = true;
    }
    return ok;
}
\end{lstlisting}
\textbf{Engineering reasoning.} A bounds argument plus validation keeps access within the declared valid extent if the caller provides an accurate buffer length. The contract still requires proof that n is valid for data.

\textbf{Verification exercise.} Try n=0, index=n-1 and index=n; add static assumptions for pointer provenance.

\subsection{Order-dependent expression effects}
\TrainingTag{Cross-reference: Rule 13.2; Rule 13.3}
\subsubsection*{Problematic}
\begin{lstlisting}[style=engineerc]
uint32_t i = 1U;
uint32_t x = i++ + i; /* unsequenced modification/read */
\end{lstlisting}
\subsubsection*{Preferred pattern}
\begin{lstlisting}[style=engineerc]
uint32_t i = 1U;
uint32_t before = i;
i = i + 1U;
uint32_t x = before + i;
\end{lstlisting}
\textbf{Engineering reasoning.} Separating state changes from value calculations makes sequencing and intent explicit. Check overflow bounds on the increment separately.

\textbf{Verification exercise.} Review each increment or decrement for unrelated reads or effects in the same full expression.

\subsection{Logical short-circuit with side effect}
\TrainingTag{Cross-reference: Rule 13.5}
\subsubsection*{Problematic}
\begin{lstlisting}[style=engineerc]
bool permit(bool ready)
{
    return ready && clear_fault_flag(); /* mutates state */
}
\end{lstlisting}
\subsubsection*{Preferred pattern}
\begin{lstlisting}[style=engineerc]
bool permit(bool ready)
{
    bool cleared = false;
    if (ready)
    {
        cleared = clear_fault_flag();
    }
    return cleared;
}
\end{lstlisting}
\textbf{Engineering reasoning.} Mutation on the right side of a short-circuit operator is hard to reason about. Explicit control flow expresses when the mutation is intended.

\textbf{Verification exercise.} Add a test that the fault-clear function is not invoked when ready is false.

\subsection{Boolean control expression}
\TrainingTag{Cross-reference: Rule 14.4}
\subsubsection*{Problematic}
\begin{lstlisting}[style=engineerc]
uint32_t alarms = get_alarm_count();
if (alarms) { activate_buzzer(); }
\end{lstlisting}
\subsubsection*{Preferred pattern}
\begin{lstlisting}[style=engineerc]
uint32_t alarms = get_alarm_count();
if (alarms != 0U) { activate_buzzer(); }
\end{lstlisting}
\textbf{Engineering reasoning.} Make the predicate evaluate to an essentially Boolean value rather than relying on integer truthiness.

\textbf{Verification exercise.} Flag bare integer predicates and ensure comparison types agree.

\subsection{Switch fallback}
\TrainingTag{Cross-reference: Rule 16.4; Rule 16.3}
\subsubsection*{Problematic}
\begin{lstlisting}[style=engineerc]
switch (mode)
{
    case MODE_OFF: stop_motor(); break;
    case MODE_ON:  start_motor(); break;
}
\end{lstlisting}
\subsubsection*{Preferred pattern}
\begin{lstlisting}[style=engineerc]
switch (mode)
{
    case MODE_OFF: stop_motor(); break;
    case MODE_ON:  start_motor(); break;
    default:       enter_safe_state(); break;
}
\end{lstlisting}
\textbf{Engineering reasoning.} The explicit default makes out-of-range values a deliberate safety case. Also check switch structure, clauses and type constraints.

\textbf{Verification exercise.} Test a corrupted mode value and prove the safe-state branch is reached.

\subsection{Error return handling}
\TrainingTag{Cross-reference: Dir 4.7; Rule 17.7}
\subsubsection*{Problematic}
\begin{lstlisting}[style=engineerc]
(void)driver_write(frame); /* status discarded */
\end{lstlisting}
\subsubsection*{Preferred pattern}
\begin{lstlisting}[style=engineerc]
driver_status_t status = driver_write(frame);
if (status != DRIVER_OK)
{
    record_driver_failure(status);
    enter_safe_state();
}
\end{lstlisting}
\textbf{Engineering reasoning.} Casting a meaningful status to void does not satisfy the directive to check error information. Exact handling must be established by the API contract.

\textbf{Verification exercise.} Fault-inject all documented return statuses and assert error recovery.

\subsection{Dynamic allocation policy}
\TrainingTag{Cross-reference: Dir 4.12; Rule 21.3}
\subsubsection*{Problematic}
\begin{lstlisting}[style=engineerc]
uint8_t *packet = malloc(packet_length);
/* allocation and lifetime are dynamic */
\end{lstlisting}
\subsubsection*{Preferred pattern}
\begin{lstlisting}[style=engineerc]
#define PACKET_CAPACITY (256U)
static uint8_t packet[PACKET_CAPACITY];
/* validate packet_length <= PACKET_CAPACITY before use */
\end{lstlisting}
\textbf{Engineering reasoning.} A project-controlled bounded storage allocation avoids dependence on heap exhaustion and fragmentation; static storage must still be concurrency-safe and sized correctly.

\textbf{Verification exercise.} Verify there are no calls to the prohibited library allocators and that capacity checks precede reads/writes.

\subsection{Macro parameter precedence}
\TrainingTag{Cross-reference: Rule 20.7; Dir 4.9}
\subsubsection*{Problematic}
\begin{lstlisting}[style=engineerc]
#define DOUBLE_SUM(x) 2U * x
uint32_t a = DOUBLE_SUM(3U + 1U);
\end{lstlisting}
\subsubsection*{Preferred pattern}
\begin{lstlisting}[style=engineerc]
static inline uint32_t double_value(uint32_t x)
{
    return x * 2U;
}
uint32_t a = double_value(3U + 1U);
\end{lstlisting}
\textbf{Engineering reasoning.} Functions provide type-checked argument evaluation and a simpler review surface. The expression still needs a documented numeric range.

\textbf{Verification exercise.} Expand the original macro and inspect the actual precedence; audit overflow contracts.

\subsection{Public linkage discipline}
\TrainingTag{Cross-reference: Rule 8.4; Rule 8.8}
\subsubsection*{Problematic}
\begin{lstlisting}[style=engineerc]
int32_t calibration_count = 0; /* global by default */
\end{lstlisting}
\subsubsection*{Preferred pattern}
\begin{lstlisting}[style=engineerc]
/* Internal module state */
static int32_t calibration_count = 0;
/* Public definitions belong behind an external declaration. */
\end{lstlisting}
\textbf{Engineering reasoning.} Limit linkage to the smallest necessary scope and ensure any external definition has a compatible visible declaration.

\textbf{Verification exercise.} Inspect the link map and symbol table; compare external declarations across translation units.

\subsection{String parsing and bounds}
\TrainingTag{Cross-reference: Dir 4.14; Rule 21.17}
\subsubsection*{Problematic}
\begin{lstlisting}[style=engineerc]
void receive(const char *incoming)
{
    char buffer[8];
    strcpy(buffer, incoming); /* no length proof */
}
\end{lstlisting}
\subsubsection*{Preferred pattern}
\begin{lstlisting}[style=engineerc]
bool receive(const char *incoming, size_t n)
{
    bool ok = false;
    static char buffer[8];
    if ((incoming != NULL) && (n < sizeof buffer))
    {
        for (size_t i = 0U; i < n; ++i)
        {
            buffer[i] = incoming[i];
        }
        buffer[n] = '\0';
        ok = true;
    }
    return ok;
}
\end{lstlisting}
\textbf{Engineering reasoning.} This bounded-copy illustration assumes incoming designates at least n readable bytes. The input length must come from an independently validated transport contract.

\textbf{Verification exercise.} Exercise n=0, 7, 8; fuzz length metadata and check memory safety.

\subsection{Resource lifecycle}
\TrainingTag{Cross-reference: Rule 22.1; Rule 22.6}
\subsubsection*{Problematic}
\begin{lstlisting}[style=engineerc]
FILE *f = fopen("status.log", "r");
if (f != NULL) { return 0; } /* stream leaked */
\end{lstlisting}
\subsubsection*{Preferred pattern}
\begin{lstlisting}[style=engineerc]
FILE *f = fopen("status.log", "r");
if (f != NULL)
{
    int close_status = fclose(f);
    f = NULL;
    if (close_status != 0) { /* report failure */ }
}
\end{lstlisting}
\textbf{Engineering reasoning.} For projects permitting hosted I/O, release acquired resources along all paths. The preferred fragment is for lifecycle instruction; Rule 21.6 can separately forbid standard I/O in the chosen profile.

\textbf{Verification exercise.} Inspect early returns, error paths and permitted APIs; use a resource ownership table.

\subsection{Thread-shared state}
\TrainingTag{Cross-reference: Dir 5.1; Dir 5.2}
\subsubsection*{Problematic}
\begin{lstlisting}[style=engineerc]
/* Shared across C11 threads; read and write unsynchronized */
static uint32_t total = 0U;
int worker(void *arg) { total++; return 0; }
\end{lstlisting}
\subsubsection*{Preferred pattern}
\begin{lstlisting}[style=engineerc]
#include <stdatomic.h>
static atomic_uint total = ATOMIC_VAR_INIT(0U);
int worker(void *arg)
{
    (void)arg;
    (void)atomic_fetch_add_explicit(
        &total, 1U, memory_order_seq_cst);
    return 0;
}
\end{lstlisting}
\textbf{Engineering reasoning.} The atomic demonstration addresses a race on the counter, but full C11 threading analysis must still review allowed operations, initialization, resource lifetime and sequential-consistency constraints.

\textbf{Verification exercise.} Run repeated stress tests and an approved thread checker; inspect data paths not covered by atomic accesses.

\subsection{Generic selection and side effects}
\TrainingTag{Cross-reference: Rule 23.7; Rule 23.2}
\subsubsection*{Problematic}
\begin{lstlisting}[style=engineerc]
#define SELECT_AND_DOUBLE(x) \
    _Generic((x), int: ((x) + (x)), default: 0)
int k = 1;
int result = SELECT_AND_DOUBLE(k++); /* repeated effects */
\end{lstlisting}
\subsubsection*{Preferred pattern}
\begin{lstlisting}[style=engineerc]
static int doubled_int(int x) { return x + x; }
#define SELECT_AND_DOUBLE(x) \
    _Generic((x), int: doubled_int, default: doubled_int)(x)
int k = 1;
int result = SELECT_AND_DOUBLE(k++);
\end{lstlisting}
\textbf{Engineering reasoning.} \_Generic controlling expressions are not evaluated, but selected association expressions can evaluate macro arguments; using a selected function limits the value evaluation to one call.

\textbf{Verification exercise.} Inspect preprocessed output for every supported type and count side effects.

\subsection{Shifts and operand widths}
\TrainingTag{Cross-reference: Rule 12.2}
\subsubsection*{Problematic}
\begin{lstlisting}[style=engineerc]
uint32_t shift(uint32_t v, uint32_t count)
{
    return v << count; /* count may be too large */
}
\end{lstlisting}
\subsubsection*{Preferred pattern}
\begin{lstlisting}[style=engineerc]
bool shift(uint32_t v, uint32_t count, uint32_t *out)
{
    bool ok = false;
    if ((out != NULL) && (count < 32U))
    {
        *out = v << count;
        ok = true;
    }
    return ok;
}
\end{lstlisting}
\textbf{Engineering reasoning.} The illustrative guard assumes uint32\_t has exactly 32 value bits. The target type model must support this assumption; shift-result policy may need additional overflow guards.

\textbf{Verification exercise.} Test edge shift counts and prove that the source type is exactly the expected width.


\section{Cross-cutting example review}
For every example: identify all applicable rules, examine implementation-dependent assumptions, run compile-time diagnostics and static checks, exercise boundaries in unit tests, and create a requirements link. An isolated snippet cannot prove entire-system conformance, particularly for system-scope or undecidable guidelines.

\end{document}
